% Local mathematical reconstructions of Ganeri2004 pp.309--391.
% Not proofs of historical adequacy or a complete Indian logical calculus.
\begin{annexitem}{annex:p13.ganeri.bridge}
 {thm:p13-ganeri-bridge}{Explicit bridges and insufficient resemblance}
\begin{proof}
Instantiate \(\forall x(Fx\to Gx)\) at \(a\), obtaining \(Fa\to Ga\).
Modus ponens with \(Fa\) yields \(Ga\).
For the second claim, assume \(P\). Modus ponens with \(P\to Q\) yields
\(Q\), contradicting \(\neg Q\). Discharging the assumption gives
\(\neg P\). This direction needs an explicit implication or a derivation
of \(Q\) from the other commitments and \(P\); the bare words
relevant resemblance do not supply it.

For the last claim take domain \(\{a,b\}\), with distinct named elements,
\(F=\{a,b\}\), and \(G=\{b\}\). Then \(Fa,Fb,Gb\) are true and
\(Ga\) false. This is a countermodel to their alleged consequence.
The universal bridge fails at \(a\), so the countermodel does not refute
the first claim. Thus it identifies precisely the omitted premise.
\end{proof}
\end{annexitem}

\begin{annexitem}{annex:p13.ganeri.fourfold}
 {thm:p13-ganeri-fourfold}{Partition, root cases, and equivalence counterexample}
\begin{proof}
Let \(B=R(a)\). If \(B=\varnothing\), the fourth case holds and no other
case does. Otherwise test \(a\in B\). If false, the second case holds.
If true, either every member of \(B\) equals \(a\), in which case
\(B=\{a\}\), or some \(b\in B\) differs from \(a\), in which case
the third holds. These tests are exhaustive in classical set theory.
The first and third are distinguished by the existence of a distinct
member, and both exclude the second and fourth; this proves disjointness.

For \(b^2=a\) in the reals, \(a=0\) has the singleton root set \(\{0\}\).
At \(a=4\), the roots are \(2,-2\), neither equal to \(4\).
At \(a=1\), the roots are \(1,-1\), so the root set contains \(a\)
and another member. At \(a=-1\) no real square is negative, so the
root set is empty. Finally, on \(\{a,b\}\) with \(a\ne b\), the universal
relation is reflexive, symmetric, and transitive; its class at \(a\)
is \(\{a,b\}\). It realizes the third case despite being an equivalence
relation. Identity excludes that case because its classes are singletons,
not merely because it is transitive.
\end{proof}
\end{annexitem}

\begin{annexitem}{annex:p13.ganeri.pervasion}
 {thm:p13-ganeri-pervasion}{A preorder characterization}
\begin{proof}
If \(F\preceq G\) and \(G\preceq Z\), transitivity gives
\(F\preceq Z\). This proves the right-hand universal condition.
Conversely substitute \(Z=G\) in that condition. Reflexivity gives
\(G\preceq G\), so the instantiated implication gives \(F\preceq G\).
No antisymmetry is used, and no existence of an empirical cause is inferred.
\end{proof}
\end{annexitem}

\begin{annexitem}{annex:p13.ganeri.wheel}
 {thm:p13-ganeri-wheel}{A distribution classification is not knowledge of the site}
\begin{proof}
For nonempty \(C\), either \(H\cap C=\varnothing\), or it is nonempty.
In the latter case either \(H\cap C=C\), or it is a proper subset.
These are respectively \(-,+,\pm\). Nonemptiness of \(C\) makes the
first two distinct, and a proper nonempty subset distinguishes the third.
There are therefore \(3\cdot3=9\) ordered pairs.
Positive occurrence requires row \(+\) or row \(\pm\); absence from
\(\mathsf{vp}\) requires column \(-\). Row-major numbering places
these at 2 and 8, and excludes the other seven.
Each pair is realizable: take each comparison class with two members
and choose all, none, or one of them for \(H\). This also verifies
that the classification has not excluded one of its formal cells.

To refute the alleged site consequence take
\(D=\{p,a,b\}\), \(\mathsf{sp}=\{a\}\), \(\mathsf{vp}=\{b\}\),
\(H=\{p,a\}\), and \(S=\{a\}\). The reason is attached at \(p\),
is positive on every member of \(\mathsf{sp}\), and absent from
\(\mathsf{vp}\). It passes cell 2, but \(S(p)\) is false.
The comparison classes accurately record \(S\) at their members;
they simply do not settle the disputed site.

For an empty \(C\), \(H\cap C=C\) and \(H\cap C=\varnothing\)
both hold. Thus assigning either \(+\) or \(-\) by those equations
would lose the distinction. A separate \(0\) status records the empty
range. The expanded wheel's historical acceptance choices are additional
criteria; they are not conclusions of the positive-witness test proved here.
\end{proof}
\end{annexitem}

\begin{annexitem}{annex:p13.ganeri.crosssite}
 {thm:p13-ganeri-crosssite}{Full existential witness transfer}
\begin{proof}
Eliminate \(\exists xFx\) using a fresh name \(a\) and temporary \(Fa\).
Universal instantiation of \(\forall x\exists y(Fx\to Gy)\) gives
\(\exists y(Fa\to Gy)\). Eliminate this existential with a name \(b\)
fresh for all undischarged premises and the intended conclusion.
The temporary premise \(Fa\to Gb\), with \(Fa\), gives \(Gb\).
Existential introduction gives \(\exists yGy\).
Discharge the witness assumption for \(b\), then the witness assumption
for \(a\). The conclusion contains neither fresh name, so both
existential eliminations meet their side conditions.

For the last claim take two objects \(a,b\), \(F=\{a\}\), \(G=\{b\}\).
For \(x=a\), choose \(y=b\). For \(x=b\), the implication has a false
antecedent and any \(y\) suffices. Both premises are true, while \(Ga\)
is false. The result therefore transfers existence, not identity of sites.
\end{proof}
\end{annexitem}

\begin{annexitem}{annex:p13.ganeri.standpoints}
 {thm:p13-ganeri-standpoints}{Realizing seven patterns and locating adjunction}
\begin{proof}
Fix a classical propositional language with distinct atoms \(p,q\).
For a set of formulas \(A\), let \(\operatorname{Cn}(A)\) be all its
classical semantic consequences. A standpoint is a consistent theory
\(\operatorname{Cn}(A)\); write \(\sigma:r\) for membership.
Consider \(\sigma_t=\operatorname{Cn}(\{p\})\),
\(\sigma_f=\operatorname{Cn}(\{\neg p\})\), and
\(\sigma_u=\operatorname{Cn}(\varnothing)\).

The first has valuations with \(p\) true, so is consistent. It contains
\(p\), and cannot contain \(\neg p\) because such a valuation refutes it.
The second analogously asserts \(\neg p\) but not \(p\).
In the third, a valuation with \(p=0\) refutes \(p\), and one with \(p=1\)
refutes \(\neg p\); hence it asserts neither. Every nonempty subset
of these three standpoints realizes exactly its associated subset
of \(\{t,f,u\}\). There are \(2^3-1=7\) such subsets. These examples
also show that the support markers are not exclusive over a whole discourse.

For failure of global adjunction take the two standpoints
\(\operatorname{Cn}(\{p\})\) and \(\operatorname{Cn}(\{q\})\).
The first asserts \(p\), the second \(q\). A valuation with \(p=1,q=0\)
is a model of the first but refutes \(p\land q\); a valuation with
\(p=0,q=1\) does the same for the second. Thus neither standpoint
asserts the conjunction.
Within one standpoint, if every model satisfies \(p\) and every model
satisfies \(q\), every model satisfies \(p\land q\).

This model validates the proposed kernel's nonvacuity because every
standpoint contains a tautology. Consistency excludes membership of
\(r\land\neg r\). If \(r\) entails \(s\) and a standpoint contains \(r\),
all its models satisfy \(s\), so it contains \(s\). The optional
exclusion of neutrality is not valid in this model: \(\sigma_u\)
is its explicit counterexample. No modal completeness claim is needed.
\end{proof}
\end{annexitem}

\begin{annexitem}{annex:p13.ganeri.correlation}
 {thm:p13-ganeri-correlation}{Conjunction depends on matching standpoints}
\begin{proof}
A valuation \(v\) supplies the complete standpoint consisting of all
formulas true at \(v\). In family \(A\), take valuations
\((p,q)=(1,1),(0,0)\). Each atom is true at one standpoint and false
at the other, so each has the mixed pattern. Their conjunction is
true at the first and false at the second, also mixed.
In family \(B\), take valuations \((1,0),(0,1)\). The atom patterns
are unchanged, but \(p\land q\) is false at both standpoints.
Each standpoint is consistent and complete because it comes from
a classical valuation. The two families therefore meet identical
marginal assumptions but give different conjunction patterns.
This proves the failure of a single-valued table on those marginals;
it does not prohibit a richer semantics that records their correlations.
\end{proof}
\end{annexitem}

\begin{annexitem}{annex:p13.ganeri.graph}
 {thm:p13-ganeri-graph}{Finite levels and necessary additional hypotheses}
\begin{proof}
Use exactly the node and edge constraints of the presentation.
A cycle cannot contain a universal or individuator because such nodes
have indegree zero. It cannot contain a quality or motion because every
predecessor of such a node is a universal or individuator, already
excluded from cycles. A remaining cycle would consist entirely of
substances. Transitivity of the strict proper-part relation along its
edges would give a substance properly part of itself, contradicting
irreflexivity. Thus the graph is acyclic.

In a finite acyclic graph a directed path cannot repeat a vertex:
a repeated vertex would close a cycle. Its length is therefore bounded
by the number of vertices minus one. Every vertex has a length-zero
incoming path, so the set of incoming path lengths is nonempty, finite
and bounded, and has a maximum.
Universals and individuators have no incoming edges, hence level zero.
Every quality or motion has a predecessor of level zero, and all
its predecessors are of those types, so its longest incoming path
has length one. Every substance is reached from a quality or motion,
which itself has such a predecessor. It therefore has an incoming
path of length at least two. These prove all stated levels.

For the strict-part condition's necessity, take a universal \(u\),
a quality \(q\), substances \(s_0,s_1,s_2\), and edges
\(u\to q\to s_0\to s_1\to s_2\to s_0\).
Every substance is reached from \(q\), and the type conditions hold,
but the substance cycle cannot belong to a strict transitive relation.
Without that condition the other constraints permit it.
For the finiteness condition's necessity take \(u,q,s_0,s_1,\ldots\),
edges \(u\to q\), \(q\to s_n\) for every \(n\), and
\(s_{n+1}\to s_n\) for every \(n\). On the substances use the
strict transitive relation \(s_m<s_n\) exactly when \(m>n\).
There is no cycle and all type and reachability conditions hold.
For each \(n\), the path
\(u,q,s_n,s_{n-1},\ldots,s_0\) has length \(n+2\).
Thus \(s_0\) has no longest finite incoming path. This refutes the
unqualified inference from acyclicity to the existence of a level.
\end{proof}
\end{annexitem}

\begin{annexitem}{annex:p13.ganeri.absence}
 {thm:p13-ganeri-absence}{Higher negations and the difference between two consequences}
\begin{proof}
Write \(b_n=T(\neg^np)\). At depth zero the allowed pairs
\((b_0,b_1)\) exclude \((0,0)\). Hence \(b_0=0\) implies \(b_1=1\).
For every depth \(n\geq1\), the recurrence gives
\(b_{n+1}=1-b_n\), so \(b_n=0\) implies \(b_{n+1}=1\).
This proves the first constraint at every depth.
The second constraint at depth \(n\) is
\(b_{n+2}=1-b_{n+1}\), exactly an instance of the recurrence.

If \(b_2=1\), then \(b_1=0\); exclusion of \((0,0)\) forces \(b_0=1\).
This proves double-negation elimination at depth zero.
At every later depth the recurrence is periodic with period two,
so \(b_{n+2}=b_n\), proving elimination there as well.
In the allowed initial case \((b_0,b_1)=(1,1)\), one has \(b_2=0\).
Thus \(p\) is true while its double negation is false, refuting
double-negation introduction. The same case refutes
\(T\neg p\Rightarrow\neg Tp\).

For the two consequence conditions, a one-location example suffices:
give \(\alpha\) initial pair \((1,0)\) and \(\beta\) pair \((1,1)\).
Both satisfy the constraints by the recurrence just proved.
Truth of \(\alpha\) preserves truth of \(\beta\) at every location
(there is only one). But \(\neg\beta\) is true there and
\(\neg\alpha\) is not. Preservation of object-level absence in the
contraposed direction therefore fails. This separates the conditions
without assuming a semantics for arbitrary binary connectives.
\end{proof}
\end{annexitem}

\begin{annexitem}{annex:p13.ganeri.number}
 {thm:p13-ganeri-number}{Unique possession and collective counting}
\begin{proof}
Suppose \(Fx\) and \(\neg\exists y(Fy\land y\ne x)\).
For arbitrary \(a\), assume \(Fa\). If \(a\ne x\), conjunction and
existential introduction would yield the prohibited existential.
Classical reasoning therefore gives \(a=x\); discharge \(Fa\)
and universally generalize. The name \(a\) was arbitrary and fresh
for the undischarged premises. This proves
\(Fx\land\forall y(Fy\to y=x)\).
Conversely suppose this latter formula. A witness to
\(\exists y(Fy\land y\ne x)\) would, by universal instantiation
and modus ponens, be equal to \(x\), contradicting its inequality.
Discharge the witness to obtain the negated existential and retain \(Fx\).
Both directions are proved; the classical step in the first direction
is stated rather than attributed to an unrestricted historical calculus.

For the collective model take distinct pots \(a,b\), the finite set
\(C=\{a,b\}\), and a number node for 2. A collective edge from that
node to \(C\) is true by the stipulated condition \(|C|=2\).
Neither singleton \(\{a\}\) nor \(\{b\}\) has cardinality 2.
Thus no such edge to either singleton follows. If individual numerical
possession is interpreted by its singleton, each pot is one, while
their collection is two. This is a countermodel to the proposed
distributive consequence, not an existence theorem for unrestricted
collective objects.
\end{proof}
\end{annexitem}
